% Isolated post-freeze diagnostic candidate; no canonical assertion yet.
\subsection{An exceptional wall inside the prime-five Heisenberg family}

Let \(k\) be algebraically closed of characteristic two, and take
the prime-five matrices \(B_y,M_s\) of (PHR2), with
\(c_1=1\) and \(c_2=t\). This is an abstract Heisenberg tensor;
no Cartier realization on an abelian surface is assumed.

\begin{lemma}[Exact adjugate and inverse-composition identities]
\label{lem:prime-five-exceptional-composition}
\status{Candidate; exact polynomial proof awaits independent audit.}
Let \(\sigma(y)\) be the five signed maximal Pfaffians of \(B_y\)
(signs are immaterial in characteristic two). There is a homogeneous
cubic vector \(\tau(s)\), defined coefficientwise over
\(\mathbf F_2[c_1,c_2]\), with
\[
       \operatorname{adj}(M_s)=s\tau(s)^t,
       \qquad \tau(\sigma(y))=E(y)y,                 \tag{PEW1}
\]
where \(E\) is homogeneous of degree five in \(y\).
On every specialization \(c_1=1,c_2=t\), the polynomial \(E\)
is nonzero. Put
\[
                         Q(y)=\sum_i y_i\sigma_i(y)^2.
                                                        \tag{PEW2}
\]
Then \(E\) is a nonzero scalar multiple of \(Q\) if and only if
\[
                              t^5+1=0.               \tag{PEW3}
\]
\end{lemma}
\begin{proof}
The retained standalone program expands all five Pfaffians and
all twenty-five four-by-four cofactors over the binary polynomial
ring in \(c_1,c_2,y_0,\ldots,y_4\). Dividing the first adjugate
row by the corresponding section coordinate defines \(\tau\);
the program verifies every entry of the first identity (PEW1).
It substitutes \(\sigma\) into every entry of \(\tau\), divides
the \(j\)-th entry by \(y_j\), and verifies that all five
quotients agree. This proves the second identity by exact
polynomial reconstruction, including all zero-coordinate cases.

Before specialization, \(Q\) has fifteen binary monomials and
\(E\) has thirty-seven. The output retains every exponent and
coefficient. Write \(Q=\sum q_\nu(c)y^\nu\) and
\(E=\sum e_\nu(c)y^\nu\). Choose the retained index \(\nu_0\)
for which \(q_{\nu_0}=c_1^4\). On setting \(c_1=1,c_2=t\),
the program computes in \(\mathbf F_2[t]\) the exact identities
\[
 \gcd_\nu e_\nu(1,t)=1,\qquad
 \gcd_\nu\bigl(e_\nu(1,t)+q_\nu(1,t)e_{\nu_0}(1,t)\bigr)
                              =t^5+1.                \tag{PEW4}
\]
The first ensures that \(E\) is nonzero at every parameter.
Because \(q_{\nu_0}(1,t)=1\), the second says exactly that
all coefficients of \(E-e_{\nu_0}Q\) vanish precisely at
(PEW3). At such a parameter the scalar \(e_{\nu_0}\) cannot
be zero, by the first identity. This proves the equivalence.
These are polynomial gcd calculations, not testing finitely many
parameter values.
\end{proof}

\begin{proposition}[The common maximal-rank Euler wall can be empty]
\label{prop:prime-five-empty-common-wall}
\status{Candidate; complete abstract consequence awaits independent audit.}
For \(t^5=1\), the abstract tensor above has generic dual and
multiplication ranks four and a birational radical map, but there
is no section \([s]\) with all three properties:
\[
 \operatorname{rank}M_s=4,\qquad
 \operatorname{rank}B_{y_s}=4,\qquad
 y_s(s^{[2]})=0,
 \quad\text{where }\ker M_s^t=k y_s.                 \tag{PEW5}
\]
Thus the two generic rank statements and birationality do not by
themselves prove a nonempty common maximal-rank Euler wall, even
within the exact prime-dimensional Heisenberg tensor family.
\end{proposition}
\begin{proof}
The generic ranks and birationality follow from the canonical
prime-five theorem. On the dual maximal-rank open,
\(s=\sigma(y)\ne0\). If \(Q(y)=0\), then (PEW3) and
(PEW1) give \(\tau(s)=0\), so
\(\operatorname{adj}(M_s)=0\) and \(\operatorname{rank}M_s<4\).

Conversely, if all conditions (PEW5) held, the relation
\(M_s^t y_s=0\) would give \(B_{y_s}s=0\). Dual maximal
rank would force \([s]=[\sigma(y_s)]\). The third condition
would imply \(Q(y_s)=0\), contradicting the preceding paragraph.
This proves the asserted emptiness without any specialization
across a rank drop.
\end{proof}

\begin{remark}[Actual Cartier realization remains a separate question]
\label{rem:prime-five-empty-wall-scope}
\status{Candidate scope; no geometric realization claimed.}
The exact executed program is
\path{code/monodromy/post_CP36_prime5_wall_diagnostic.py}; it passed
at 05:44:26.793536 UTC on 20 November 2026 in
\(0.003974181992816739\) seconds. Its complete polynomial vectors,
all proportionality coefficient equations and the two exact gcds
are retained in the adjacent author output. No execution failed.
This is an obstruction to an algebraic inference, not a
counterexample to the generic curve-rank theorem or to any full
monodromy prediction. No assertion is made that the exceptional
coefficient ratios occur for an actual globally generated
type-\((1,5)\) Cartier tensor. The current checkpoint inputs are
unchanged and this note is unpromoted.
\end{remark}
